\chapter{Vitrectomy}

\section*{What Is This Surgery?}
Vitrectomy is a specialized surgical procedure performed by vitreoretinal surgeons to excise the vitreous gel, a transparent, jelly-like substance filling the posterior segment of the eye. This surgery is essential for addressing various severe vitreoretinal disorders that pose a significant risk to vision, including retinal detachment, vitreous hemorrhage, macular holes, and epiretinal membranes. By removing the vitreous, surgeons gain direct access to the retina, enabling them to repair underlying pathologies, relieve traction, and clear opacities. Vitrectomy is often a vision-saving intervention, preventing permanent blindness or severe visual impairment associated with these conditions. The procedure can be performed using various techniques, including traditional and minimally invasive approaches, depending on the specific clinical scenario.

\section*{Alternative Names}
\begin{itemize}
  \item Pars plana vitrectomy (PPV)
  \item Posterior vitrectomy
  \item Vitreous removal surgery
  \item Vitreoretinal surgery
  \item Endovitrectomy
  \item Vitrectomy with membrane peeling
\end{itemize}

\section*{Relevant Surgical Anatomy}
The surgical anatomy relevant to vitrectomy includes the vitreous body, retina, choroid, and surrounding structures. The vitreous gel is a clear, gelatinous substance that occupies the space between the lens and the retina, providing structural support to the eye. 

The retina is a thin layer of neural tissue that lines the back of the eye and is responsible for converting light into neural signals. It consists of several layers, including the photoreceptor layer, which contains rods and cones, and the retinal pigment epithelium (RPE), which supports photoreceptor function. 

Key vascular supply to the retina is provided by the central retinal artery, a branch of the ophthalmic artery, while venous drainage occurs through the central retinal vein. The choroid, located between the retina and sclera, is rich in blood vessels and provides nutrients to the outer layers of the retina.

Important surgical landmarks include:

\begin{itemize}
  \item **Pars Plana:** The region of the ciliary body where sclerotomies are typically made for instrument access.
  \item **Macula:** The central part of the retina responsible for high-acuity vision.
  \item **Optic Nerve Head:** The point where the optic nerve exits the eye, crucial for assessing retinal health.
  \item **Calot's Triangle:** An anatomical landmark used in various ocular surgeries, although more relevant in cholecystectomy, it emphasizes the importance of understanding surrounding structures.
\end{itemize}

Understanding these anatomical relationships is critical for successful vitrectomy, as it allows the surgeon to navigate the delicate structures of the eye while minimizing the risk of complications.

\section*{Surgical Principle \& Rationale}
The primary principle of vitrectomy is the surgical excision of the vitreous humor to eliminate pathological conditions within the vitreous cavity or to facilitate access to the retina and choroid for repair. The vitreous gel, while normally transparent, can become opaque due to hemorrhage, inflammation, or infection, or it can exert abnormal traction on the retina, leading to tears, detachments, or macular distortion. By excising the vitreous, the surgeon can directly address these issues.

The procedure employs specialized micro-instruments inserted through small incisions (sclerotomies) in the pars plana, a non-seeing part of the eye. A high-speed vitrector cuts and aspirates the vitreous gel, while an infusion line maintains intraocular pressure and volume. Illumination is provided by a fiber optic light probe. Once the vitreous is removed, the surgeon can perform various maneuvers such as peeling epiretinal membranes or the internal limiting membrane, draining subretinal fluid, performing endolaser photocoagulation to seal retinal breaks, or removing foreign bodies. The vitreous cavity is then typically filled with a balanced salt solution, a gas bubble (e.g., SF6, C3F8), or silicone oil, which serves as a temporary or long-term tamponade to support the retina while it heals. The choice of tamponade agent depends on the specific pathology and the need for prolonged internal support.

\section*{History \& Surgical Evolution}
The evolution of vitrectomy as a surgical technique began in the mid-20th century, with significant milestones marking its development. In 1971, Dr. Robert Machemer performed the first successful pars plana vitrectomy at the Bascom Palmer Eye Institute, revolutionizing the approach to vitreoretinal surgery. His innovation involved the creation of a single-port instrument capable of cutting, aspirating, and infusing, which allowed for the effective removal of opaque vitreous and the repair of complex retinal detachments.

Following Machemer's pioneering work, the technique underwent rapid refinement. In 1974, Dr. Conor O'Malley and Dr. Ralph Heinz introduced the three-port system, which became the standard for vitrectomy. This system allowed for continuous infusion, illumination, and vitreous removal through separate ports, greatly enhancing surgical control and safety. Over the subsequent decades, advancements in instrumentation led to the miniaturization of tools from 20-gauge to 23-, 25-, and even 27-gauge systems, resulting in smaller, often sutureless incisions and faster recovery times. The introduction of wide-angle viewing systems, high-speed vitrectors, and sophisticated endolaser and membrane peeling instruments has further expanded the indications and improved outcomes of vitrectomy, solidifying its role as a cornerstone of modern vitreoretinal surgery.

\section*{Clinical Indications}
Vitrectomy is indicated for a variety of sight-threatening conditions affecting the vitreous and retina. Specific clinical indications include:

\begin{itemize}
  \item **Retinal Detachment:** Particularly complex rhegmatogenous detachments, tractional detachments (often associated with proliferative diabetic retinopathy), or combined tractional-rhegmatogenous detachments.
  \item **Vitreous Hemorrhage:** Non-clearing hemorrhage from conditions such as proliferative diabetic retinopathy, retinal tears, trauma, or retinal vein occlusion, which obscures vision and prevents retinal examination.
  \item **Macular Hole:** A full-thickness defect in the center of the macula, causing central vision loss and distortion.
  \item **Epiretinal Membrane (ERM):** A thin, fibrous layer that forms on the surface of the retina, causing retinal wrinkling, distortion, and decreased vision.
  \item **Diabetic Tractional Retinopathy:** Advanced diabetic eye disease where fibrous tissue grows on the retina, pulling on it and causing detachment or hemorrhage.
  \item **Endophthalmitis:** Severe intraocular infection that requires removal of infected vitreous to reduce bacterial load and allow antibiotic penetration.
  \item **Retained Lens Fragments:** Lens material remaining in the vitreous cavity after complicated cataract surgery, which can cause inflammation and elevated intraocular pressure.
  \item **Intraocular Foreign Body:** Removal of foreign objects that have penetrated the eye, often after trauma, to prevent infection and further damage.
  \item **Severe Uveitis:** Chronic or severe inflammation within the eye that may require vitreous biopsy or removal of inflammatory debris.
\end{itemize}

\section*{Clinical Positioning}
Vitrectomy can serve as both a primary and secondary surgical intervention depending on the specific vitreoretinal pathology. For many conditions, such as complex retinal detachments, non-clearing vitreous hemorrhages, and symptomatic macular holes or epiretinal membranes, vitrectomy is considered the primary and definitive surgical treatment. In these scenarios, it directly addresses the underlying anatomical or pathological issue to restore vision or prevent further deterioration. However, vitrectomy can also be a secondary or salvage procedure. This applies when initial, less invasive treatments have failed, or when complications arise from other ocular surgeries. For instance, it may be performed as a secondary procedure for recurrent retinal detachment after prior repair, to remove retained lens fragments following complicated cataract surgery, or to manage chronic endophthalmitis that has not responded to antibiotic therapy alone. Its versatility allows it to be a crucial tool at various stages of disease management.

\section*{Surgical Technique \& Procedure}
Vitrectomy is a highly specialized microsurgical procedure typically performed under local anesthesia with sedation, although general anesthesia may be used for complex cases, pediatric patients, or anxious individuals. The patient is positioned supine, and the eye is prepped and draped in a sterile fashion.

The procedure generally follows these steps:

\begin{itemize}
  \item **Anesthesia and Preparation:** After administration of local or general anesthesia, the eye is stabilized with a speculum. The surgical field is sterilized, and a drape is applied.
  
  \item **Sclerotomy Creation:** Three small incisions, usually 23- or 25-gauge (0.6-0.5 mm), are made through the pars plana, approximately 3.5-4 mm posterior to the limbus. These "ports" allow the insertion of micro-instruments. One port is for an infusion cannula to maintain intraocular pressure, another for a fiber optic light source, and the third for the vitrector and other instruments.
  
  \item **Vitreous Removal:** The vitrector, a high-speed cutting and aspiration device, is inserted into the vitreous cavity. The surgeon systematically removes the vitreous gel, often starting centrally and moving peripherally, while carefully avoiding traction on the retina.
  
  \item **Membrane Peeling and Retinal Repair:** Once the vitreous is cleared, the surgeon can address the underlying pathology. This may involve using micro-forceps to peel epiretinal membranes (ERM) or the internal limiting membrane (ILM) from the retinal surface, removing subretinal fluid, or flattening a detached retina.
  
  \item **Endolaser Photocoagulation:** If retinal tears or areas of ischemia are present, an endolaser probe is used to apply laser spots to create chorioretinal adhesions, sealing tears and preventing further detachment.
  
  \item **Fluid-Air Exchange and Tamponade:** A fluid-air exchange may be performed to remove residual fluid from the vitreous cavity. The eye is then filled with a tamponade agent, such as a gas bubble (sulfur hexafluoride or perfluoropropane) or silicone oil, to provide internal support to the retina as it heals.
  
  \item **Closure:** The micro-incisions are often self-sealing (sutureless) with smaller gauge vitrectomy. In some cases, especially with larger gauge instruments or if leakage is noted, a fine suture may be used to close the sclerotomies. Antibiotic and steroid injections may be given subconjunctivally, and an eye patch or shield is applied.
\end{itemize}

\section*{Preoperative Workup \& Preparation}
Thorough preoperative preparation is essential to ensure patient safety and optimize surgical outcomes. This typically involves a comprehensive ophthalmic and systemic evaluation.

Key preparatory steps include:

\begin{itemize}
  \item **Detailed Ophthalmic Examination:** This includes best-corrected visual acuity, slit-lamp examination, intraocular pressure measurement, and a dilated fundus examination. Optical coherence tomography (OCT) and B-scan ultrasonography may be performed to assess the retinal pathology and vitreous status, especially if media opacities preclude direct visualization.
  
  \item **Systemic Workup:** For patients undergoing general anesthesia or with significant comorbidities, a full medical evaluation is required. This may include complete blood count (CBC), electrolyte panel, coagulation studies (PT/INR, PTT), electrocardiogram (ECG), and chest X-ray.
  
  \item **Anesthesia Clearance:** Patients are evaluated by an anesthesiologist to determine fitness for anesthesia and to discuss the safest anesthetic approach.
  
  \item **Medication Review:** All current medications are reviewed. Anticoagulants or antiplatelet agents (e.g., warfarin, aspirin, clopidogrel) may need to be temporarily discontinued or adjusted in consultation with the prescribing physician to minimize bleeding risk.
  
  \item **NPO Status:** Patients are instructed to be NPO (nil per os - nothing by mouth) for a specified period (typically 6-8 hours) before surgery to prevent aspiration during anesthesia.
  
  \item **Preoperative Medications:** Topical antibiotic drops may be prescribed for several days prior to surgery to reduce the risk of infection.
  
  \item **Informed Consent:** A detailed discussion with the patient regarding the surgical procedure, its benefits, potential risks, alternatives, and expected postoperative course, including specific positioning requirements (e.g., face-down positioning for macular holes or certain retinal detachments with gas tamponade), is crucial.
\end{itemize}

\section*{Contraindications \& Precautions}
While vitrectomy is a highly effective procedure, certain conditions may contraindicate its performance or necessitate special precautions.

\textbf{Absolute Contraindications:}

\begin{itemize}
  \item Severe systemic instability that precludes safe anesthesia or surgery.
  \item Uncorrectable severe coagulopathy posing an unacceptable risk of intraoperative or postoperative hemorrhage.
  \item Active, uncontrolled ocular infection (other than endophthalmitis, which is an indication for vitrectomy) that could worsen with surgery.
  \item Complete lack of light perception in the affected eye with no potential for visual recovery.
\end{itemize}

\textbf{Relative Contraindications and Precautions:}

\begin{itemize}
  \item **Severe Corneal Opacity:** Significant corneal scarring or edema can obscure the surgeon's view, making the procedure technically challenging or impossible. Corneal transplantation may be required first.
  
  \item **Poor Visual Potential:** In cases of long-standing retinal detachment, severe optic nerve atrophy, or extensive retinal damage, the potential for visual recovery may be minimal, making the risks of surgery outweigh the benefits.
  
  \item **Advanced Glaucoma:** Patients with pre-existing severe glaucoma may be at higher risk for postoperative intraocular pressure spikes, especially with gas or silicone oil tamponade.
  
  \item **Patient's Inability to Comply with Postoperative Positioning:** Strict face-down or specific head positioning is often critical for the success of gas tamponade. Patients unable or unwilling to maintain this position may have a higher risk of failure.
  
  \item **Gas-Filled Eye and Air Travel/Altitude Changes:** Patients with intraocular gas bubbles must avoid air travel, scuba diving, or significant altitude changes until the gas has completely resorbed. Changes in atmospheric pressure can cause the gas bubble to expand, leading to dangerously high intraocular pressure and potential vision loss.
  
  \item **General Anesthesia with Nitrous Oxide:** Nitrous oxide can diffuse into intraocular gas bubbles, causing them to expand rapidly. It must be avoided during general anesthesia in patients with intraocular gas.
  
  \item **Silicone Oil Tamponade:** While effective, silicone oil may require a second surgery for removal, and it can cause complications such as emulsification, glaucoma, or keratopathy.
  
  \item **Active Ocular Inflammation:** While vitrectomy can treat some inflammatory conditions, severe active inflammation may increase surgical risks and should be managed preoperatively if possible.
\end{itemize}

\section*{Complications \& Risks}
As with any surgical procedure, vitrectomy carries potential risks and complications, which are generally discussed with the patient prior to surgery.

\begin{itemize}
  \item **General Surgical Risks:**
  
  \begin{itemize}
    \item **Anesthesia Complications:** Nausea, vomiting, allergic reactions, respiratory or cardiac events.
    \item **Bleeding:** Subconjunctival hemorrhage (bruising), vitreous hemorrhage (new or recurrent), or rare but severe choroidal hemorrhage.
    \item **Infection (Endophthalmitis):** A severe intraocular infection, though rare (0.05\%-0.1\%), can lead to devastating vision loss.
  \end{itemize}
  
  \item **Procedure-Specific Ocular Risks:**
  
  \begin{itemize}
    \item **Cataract Progression/Formation:** Vitrectomy significantly accelerates the development or progression of cataracts, especially in phakic (natural lens) eyes. Most patients will require cataract surgery within a few years.
    \item **Retinal Tear or Detachment:** A new retinal tear or detachment can occur during or after surgery, sometimes requiring further intervention.
    \item **Elevated Intraocular Pressure (IOP) / Glaucoma:** Temporary or persistent increase in IOP can occur, particularly with gas or silicone oil tamponade, potentially requiring medication or further surgery.
    \item **Hypotony:** Abnormally low intraocular pressure, which can lead to choroidal detachment or other complications.
    \item **Macular Edema:** Swelling of the macula, which can limit visual recovery.
    \item **Epiretinal Membrane (ERM) Recurrence or Formation:** New membrane growth on the retina can occur after initial removal.
    \item **Diplopia (Double Vision):** Can result from damage to extraocular muscles or changes in ocular alignment.
    \item **Corneal Edema or Damage:** Temporary or permanent corneal issues can arise.
    \item **Need for Additional Surgery:** Re-operation may be necessary for recurrent detachment, non-resolving complications, or silicone oil removal.
    \item **Vision Loss:** While the goal is to improve vision, severe complications can lead to permanent vision loss, though this is rare.
  \end{itemize}
\end{itemize}

\section*{Postoperative Recovery Timeline}
The immediate postoperative period after vitrectomy involves careful monitoring and adherence to specific instructions to ensure optimal healing and prevent complications. 

Immediately after surgery, the patient is transferred to a post-anesthesia care unit (PACU) for recovery from anesthesia. Pain medication is administered as needed, and antiemetics may be given to prevent nausea and vomiting, which can increase intraocular pressure. The eye is typically covered with a patch and a protective shield. Patients are usually discharged home the same day or the following morning. 

For patients with a gas bubble, strict head positioning (e.g., face-down, side-lying) is often crucial for several days to weeks, as instructed by the surgeon, to ensure the gas provides proper tamponade to the retina. 

Over the first 1-2 weeks, patients will use a regimen of topical eye drops, typically including antibiotics to prevent infection and corticosteroids to reduce inflammation. Vision will likely be blurry, especially if a gas bubble is present, as the gas obscures vision until it gradually resorbs (which can take 2-8 weeks depending on the type of gas). Discomfort, mild pain, and redness are common but usually subside within a few days. Activity restrictions, such as avoiding heavy lifting, bending, or strenuous exercise, are usually advised. 

Long-term, patients with silicone oil will eventually require a second surgery for its removal, typically several months after the initial vitrectomy, unless it is intended for permanent tamponade. Visual recovery is gradual and can continue for several months, depending on the underlying condition and the extent of preoperative damage.

\section*{Surgical Findings \& Pathology}
During vitrectomy, the surgeon documents various findings related to the underlying pathology. Common surgical findings include the presence of retinal tears, detachment, epiretinal membranes, or vitreous hemorrhage. The pathology report on resected tissues, such as membranes or biopsies of the vitreous, provides critical information regarding the nature of the disease process, including the presence of inflammatory cells, neovascularization, or signs of infection. This information is vital for guiding postoperative management and assessing the prognosis.

\section*{Expected Postoperative Course}
A normal, successful postoperative course following vitrectomy typically involves gradual improvement in vision and resolution of preoperative symptoms. Patients can expect some initial discomfort and blurry vision, particularly if a gas bubble is present. Over the first few weeks, vision may fluctuate as the gas bubble resorbs and the retina heals. Most patients will notice a significant reduction in symptoms such as distortion or floaters as the eye recovers. 

By the end of the first month, many patients experience improved visual acuity, although full recovery may take several months, depending on the extent of preoperative damage and the specific pathology treated. Regular follow-up visits are essential to monitor healing and address any complications that may arise.

\section*{Postoperative Care \& Follow-up}
Postoperative care is crucial for ensuring optimal recovery and minimizing complications. Patients should adhere to the following guidelines:

\begin{itemize}
  \item **Frequent Postoperative Examinations:** Initial visits are typically scheduled for day 1, week 1, and month 1 after surgery, followed by visits at 3, 6, and 12 months, or as needed. These visits involve visual acuity checks, intraocular pressure measurements, and dilated fundus examinations to assess retinal status and detect complications.
  
  \item **Suture Removal:** If sutures were used to close the sclerotomies, they may be removed at a follow-up visit, usually within 1-3 weeks.
  
  \item **Imaging:** Optical coherence tomography (OCT) may be performed to assess macular hole closure, detect recurrent epiretinal membranes, or monitor for macular edema. B-scan ultrasonography may be used if media opacities persist.
  
  \item **Refraction:** Once vision stabilizes, a new refraction for glasses may be prescribed, typically several months after surgery.
  
  \item **Cataract Surgery:** Due to the high incidence of cataract progression after vitrectomy, many patients will eventually require cataract surgery, often within 1-3 years.
  
  \item **Silicone Oil Removal:** If silicone oil was used as a tamponade, a second surgical procedure for its removal is usually planned several months after the initial vitrectomy.
  
  \item **Warning Signs:** Patients are educated to immediately contact their surgeon if they experience sudden vision loss, severe or increasing eye pain, increasing redness, discharge, new flashes of light, or an increase in floaters, as these could indicate a serious complication like infection or recurrent retinal detachment.
\end{itemize}

\section*{Epidemiology}
Vitrectomy is a frequently performed ophthalmic surgery, with its incidence steadily increasing due to an aging global population and the rising prevalence of conditions like diabetes, which often lead to vitreoretinal complications. The specific epidemiology varies by indication. For instance, proliferative diabetic retinopathy, a leading cause of vitreous hemorrhage and tractional retinal detachment, affects millions worldwide, making diabetic vitrectomy a common procedure. Rhegmatogenous retinal detachment, another major indication, has an incidence of approximately 1 in 10,000 individuals per year, with vitrectomy being the primary treatment for complex cases. Macular holes and epiretinal membranes typically affect older adults, with prevalence increasing with age. 

While vitrectomy itself has a very low mortality rate, the morbidity is primarily related to ocular complications such as cataract formation (nearly universal in phakic eyes), recurrent retinal detachment (5-15\% depending on complexity), and glaucoma (5-10\%). The procedure is performed across all age groups, from infants with retinopathy of prematurity to the elderly, with a slight male predominance in traumatic cases and a female predominance in idiopathic macular holes.

\section*{Outcomes \& Prognosis}
The outcomes and prognosis following vitrectomy are generally favorable, with high anatomical success rates for many indications, leading to significant visual improvement. For primary rhegmatogenous retinal detachment, anatomical reattachment rates typically exceed 90-95\% after a single surgery, with visual acuity improvement in a majority of cases. Macular hole closure rates are similarly high, often above 90\%, leading to improved central vision and reduced distortion. 

For vitreous hemorrhage due to proliferative diabetic retinopathy, vitrectomy effectively clears the media, and visual outcomes depend on the extent of underlying retinal damage. Prognostic factors influencing visual outcomes include the duration of symptoms before surgery, the severity and extent of the underlying pathology (e.g., presence of proliferative vitreoretinopathy in retinal detachment, size of macular hole), the patient's age, and the presence of comorbidities like advanced glaucoma. 

While anatomical success is high, functional visual recovery can be variable and may be limited by pre-existing photoreceptor damage or the development of postoperative complications such as cystoid macular edema or recurrent epiretinal membrane. Recurrence rates for retinal detachment range from 5-15\%, often requiring re-operation. Overall, vitrectomy has revolutionized the management of severe vitreoretinal diseases, offering excellent chances for vision preservation and restoration.

\section*{References}
\begin{enumerate}
  \item MedlinePlus. Vitrectomy. U.S. National Library of Medicine; 2024. Available from: \url{https://medlineplus.gov/vitrectomy.html}
  
  \item Yanoff M, Duker JS. Ophthalmology. 6th ed. Elsevier; 2024.
  
  \item American Academy of Ophthalmology. Preferred Practice Pattern: Primary Rhegmatogenous Retinal Detachment. San Francisco, CA: American Academy of Ophthalmology; 2023.
  
  \item Machemer R, Aaberg TM, Norton EWD. Vitrectomy: A Pars Plana Approach. Trans Am Acad Ophthalmol Otolaryngol. 1971;75(4):813-820.
\end{enumerate}

\clearpage